\clearpage\thispagestyle{empty}
\pdfbookmark[0]{Quellen zu den Refraktionstafeln (1822)}{refraction1822-sources}
\begin{center}\textbf{Quellen zu den Refraktionstafeln (1822)}\end{center}
\noindent Diese KI-Übertragung in heutiges Deutsch folgt den \textit{Refractionstafeln} in den \href{https://gdz.sub.uni-goettingen.de/id/PPN236007467}{\textit{Werken}, Band~VI} (Göttingen, 1874), S.~185--190. Sie enthält Heinrich Christian Schumachers Einleitung und alle drei Tafeln auf fünf Seiten.
\par\smallskip\noindent Schumacher stellt eine Umformung eines Teils der Refraktionstafel von Friedrich Wilhelm Bessel durch Carl Friedrich Gauß vor. Als Erstveröffentlichung nennt die Werkausgabe das erste Heft der von H.~C.~Schumacher herausgegebenen \textit{Sammlung von Hülfstafeln} (Kopenhagen, 1822).
\par\smallskip\noindent Übertragung und Satz wurden mit KI erstellt.
\clearpage
